%%%%%%%%%%%%%%%%%%%%%%%%%%%%%%%%%%%%%%%%%%%%%%%%%%%%%%%%%%%%%%%%%%%%%%%%%%%%%%%%
% 
% UNEMAT-BBG-Tex (v1.1)
% https://github.com/lkaranl/PRE-PROJETO-TCC-LATEX
%
% MODELO DO PRÉ-PROJETO DE TCC DA UNEMAT
% 
%%%%%%%%%%%%%%%%%%%%%%%%%%%%%%%%%%%%%%%%%%%%%%%%%%%%%%%%%%%%%%%%%%%%%%%%%%%%%%%%
% 
% MIT License
% Copyright (c) 2016 Renato Uchôa Brandão 
%
% Copyright (c) 2019 Karan Luciano Silva
% karanluciano1@gmail.com - UNEMAT BBG
% 
% Permission is hereby granted, free of charge, to any person obtaining a copy
% of this software and associated documentation files (the "Software"), to deal
% in the Software without restriction, including without limitation the rights
% to use, copy, modify, merge, publish, distribute, sublicense, and/or sell
% copies of the Software, and to permit persons to whom the Software is
% furnished to do so, subject to the following conditions:
% 
% The above copyright notice and this permission notice shall be included in all
% copies or substantial portions of the Software.
% 
% THE SOFTWARE IS PROVIDED "AS IS", WITHOUT WARRANTY OF ANY KIND, EXPRESS OR
% IMPLIED, INCLUDING BUT NOT LIMITED TO THE WARRANTIES OF MERCHANTABILITY,
% FITNESS FOR A PARTICULAR PURPOSE AND NONINFRINGEMENT. IN NO EVENT SHALL THE
% AUTHORS OR COPYRIGHT HOLDERS BE LIABLE FOR ANY CLAIM, DAMAGES OR OTHER
% LIABILITY, WHETHER IN AN ACTION OF CONTRACT, TORT OR OTHERWISE, ARISING FROM,
% OUT OF OR IN CONNECTION WITH THE SOFTWARE OR THE USE OR OTHER DEALINGS IN THE
% SOFTWARE.
%
%%%%%%%%%%%%%%%%%%%%%%%%%%%%%%%%%%%%%%%%%%%%%%%%%%%%%%%%%%%%%%%%%%%%%%%%%%%%%%%%

\documentclass{unemat-tex}

%%%%%%%%%%%%%%%%%%%%%%%%%%%%%%%%%%%%%%%%%%%%%%%%%%%%%%%%%%%%%%%%%%%%%%%%%%%%%%%%
% Inclusão de pacotes
%%%%%%%%%%%%%%%%%%%%%%%%%%%%%%%%%%%%%%%%%%%%%%%%%%%%%%%%%%%%%%%%%%%%%%%%%%%%%%%%

\usepackage{lmodern}			% Usa a fonte Latin Modern
\usepackage[T1]{fontenc}		% Selecao de codigos de fonte.
\usepackage[utf8]{inputenc}		% Codificacao do documento (conversão automática dos acentos)
\usepackage{indentfirst}		% Indenta o primeiro parágrafo de cada seção.
\usepackage{color}				% Controle das cores
\usepackage{graphicx}			% Inclusão de gráficos
\usepackage{microtype} 			% para melhorias de justificação
\usepackage{float}

% Uso da fonte Arial (UNEMAT)
\usepackage{helvet}
\renewcommand{\familydefault}{\sfdefault}

% Incluir pacotes adicionais, caso sejam necessários ....

%%%%%%%%%%%%%%%%%%%%%%%%%%%%%%%%%%%%%%%%%%%%%%%%%%%%%%%%%%%%%%%%%%%%%%%%%%%%%%%%
% Configuração das citações
%%%%%%%%%%%%%%%%%%%%%%%%%%%%%%%%%%%%%%%%%%%%%%%%%%%%%%%%%%%%%%%%%%%%%%%%%%%%%%%%
\usepackage[brazilian,hyperpageref]{backref}	 % Paginas com as citações na bibl
\usepackage[alf,
versalete,
abnt-emphasize = bf, % destaca o titulo em negrito;
abnt-etal-list = 3, % trabalhos com mais de 3 autores recebem et al.,;
abnt-etal-text = it, % escreve o et al., em italico;
abnt-and-type = &, % usa o carater '&' no lugar de 'e' para mais de um autor;
abnt-last-names = abnt, % trata sobrenomes 'estritamente' conforme a ABNT; e
abnt-repeated-author-omit = yes % autores com + de uma entrada recebem '____.'
]{abntex2cite}

% Configuração das referências bibliográficas
\renewcommand{\backref}{}
\renewcommand*{\backrefalt}[4]{	}

%%%%%%%%%%%%%%%%%%%%%%%%%%%%%%%%%%%%%%%%%%%%%%%%%%%%%%%%%%%%%%%%%%%%%%%%%%%%%%%%
% Informações da capa e da folha de rosto
%%%%%%%%%%%%%%%%%%%%%%%%%%%%%%%%%%%%%%%%%%%%%%%%%%%%%%%%%%%%%%%%%%%%%%%%%%%%%%%%

\titulo{TÍTULO DO TRABALHO: Subtítulo do trabalho}

\autor{Autor(a) do Trabalho}

\local{BARRA DO BUGRES - MT}

\data{2019}

% Alterar o nome do campus e do curso, caso houver necessidade
%\instituicao{
%	UNIVERSIDADE DO ESTADO DE MATO GROSSO\\
%	CAMPUS UNIVERSITARIO DEP. EST. RENE BARBOUR \\
%	FACULDADE DE CIENCIAS EXATTTAS E TECNOLOGICAS
%	CURSO DE CIENCIA DA COMPUTACAO
%}

\tipotrabalho{TCC}

% Informar {Titulação}{Nome}
\orientador{Prof(a). Dr(a).}{Nome do(a) Professor(a) Orientador(a)}

% Alterar o preâmbulo conforme necessário
\preambulo{Trabalho de conclusão de curso
	apresentado ao curso de Ciência da
	Computação da Universidade do Estado de
	Mato Grosso – UNEMAT, como requisito
	parcial para obtenção do grau de Bacharel.
}
% ---

% ---
% Configurações de aparência do PDF final


%%%%%%%%%%%%%%%%%%%%%%%%%%%%%%%%%%%%%%%%%%%%%%%%%%%%%%%%%%%%%%%%%%%%%%%%%%%%%%%%
% Outras configurações
%%%%%%%%%%%%%%%%%%%%%%%%%%%%%%%%%%%%%%%%%%%%%%%%%%%%%%%%%%%%%%%%%%%%%%%%%%%%%%%%

% Configuração da geração do PDF
\makeatletter
\hypersetup{
     	%pagebackref=true,
		pdftitle={\@title}, 
		pdfauthor={\@author},
    	pdfsubject={\imprimirpreambulo},
	    pdfcreator={LaTeX with abnTeX2},
		pdfkeywords={abnt}{latex}{abntex}{abntex2}{projeto de pesquisa}, 
		colorlinks=false,
		bookmarksdepth=4,
		pdfborder={0 0 0},
}
\makeatother


% O tamanho do parágrafo é dado por:
\setlength{\parindent}{1.3cm}

% Controle do espaçamento entre um parágrafo e outro:
\setlength{\parskip}{0.2cm}  % tente também \onelineskip

% compila o indice
\makeindex


%%%%%%%%%%%%%%%%%%%%%%%%%%%%%%%%%%%%%%%%%%%%%%%%%%%%%%%%%%%%%%%%%%%%%%%%%%%%%%%%
% CORPO DO TRABALHO ...
%%%%%%%%%%%%%%%%%%%%%%%%%%%%%%%%%%%%%%%%%%%%%%%%%%%%%%%%%%%%%%%%%%%%%%%%%%%%%%%%
\begin{document}

\selectlanguage{brazil}

% Retira espaço extra obsoleto entre as frases.
\frenchspacing 

% Inicializa a parte pre-textual
\pretextual

% Imprime a capa
\imprimircapa

% Imprime a folha de rosto
\imprimirfolhaderosto


% ------------------------------------------------------------------------------
% FOLHA DE APROVACAO
% ------------------------------------------------------------------------------
% Atenção: Alterar apenas a data e o nome dos convidados
\begin{folhadeaprovacao}
	
	\begin{center}
		{\ABNTEXchapterfont\bfseries\normalsize\imprimirautor}
		
		\vspace*{\fill}\vspace*{\fill}
		\begin{center}
			\ABNTEXchapterfont\bfseries\normalsize\imprimirtitulo
		\end{center}
		\vspace*{\fill}
		
		\hspace{.45\textwidth}
		\begin{minipage}{.5\textwidth}
			\imprimirpreambulo
		\end{minipage}%
		\vspace*{\fill}
	\end{center}
	
	Trabalho aprovado. \imprimirlocal, 24 de novembro de 2016:
	
	\assinatura{\textbf{\imprimirorientador} \\ Orientador} 
	\assinatura{\textbf{Prof. Dr. Fulano de tal} \\ Convidado 1}
	\assinatura{\textbf{Prof. Dr. Nome do outro convidado} \\ Convidado 2}
	
	\begin{center}
		\vspace*{0.5cm}
		{\normalsize\bfseries\imprimirlocal}
		\par
		{\normalsize\bfseries\imprimirdata}
		\vspace*{1cm}
	\end{center}
	
\end{folhadeaprovacao}


% ------------------------------------------------------------------------------
% AGRADECIMENTOS
% ------------------------------------------------------------------------------
\begin{agradecimentos}
	Folha em que o autor faz agradecimentos àqueles que contribuíram de maneira relevante para a realização do TCC...
\end{agradecimentos}


% ------------------------------------------------------------------------------
% EPIGRAFE
% ------------------------------------------------------------------------------
\begin{epigrafe}
	\vspace*{\fill}
	\begin{flushright}
		\textit{
			``Folha em que o autor apresenta uma citação,\\
			seguida da indicação da autoria, relacionada com a\\
			matéria tratada no corpo do trabalho.``\\
			(Nome do autor)
		}
	\end{flushright}
\end{epigrafe}


% ------------------------------------------------------------------------------
% RESUMO e ABSTRACT
% ------------------------------------------------------------------------------
\setlength{\absparsep}{18pt} % ajusta o espaçamento dos parágrafos do resumo

% Resumo em português
\begin{resumo}
	Apresentação concisa dos pontos relevantes do TCC, fornecendo uma visão rápida e clara dos objetivos, dos procedimentos metodológicos e das conclusões do trabalho.
	
	\textbf{Palavras-chave}: pc1, pc2, pc3 ...
\end{resumo}

% resumo em inglês (Abstract)
\begin{resumo}[Abstract]
	\begin{otherlanguage*}{english}
		English abstract version.
		
		\noindent 
		\textbf{Keywords}: kw1. kw2. kw ...
	\end{otherlanguage*}
\end{resumo}


% ------------------------------------------------------------------------------
% LISTA DE FIGURAS (Não altere nada aqui)
% ------------------------------------------------------------------------------
\pdfbookmark[0]{\listfigurename}{lof}
\listoffigures*
\cleardoublepage


% ------------------------------------------------------------------------------
% LISTA DE TABELAS (Não altere nada aqui)
% ------------------------------------------------------------------------------
\pdfbookmark[0]{\contentsname}{lot}
\listoftables*
\cleardoublepage

% ------------------------------------------------------------------------------
% LISTA DE SIGLAS E ABREVIATURAS
% ------------------------------------------------------------------------------
% Edite a lista de siglas conforme o modelo abaixo
\begin{siglas}
	\item[IBGE]{Instituto Brasileiro de Geografia e Estatística}
	\item[UNEMAT]{Universidade do Estado de Mato Grosso}
	% Incluir as siglas aqui ...
\end{siglas}


% ------------------------------------------------------------------------------
% SUMÁRIO (Não altere nada aqui)
% ------------------------------------------------------------------------------
\pdfbookmark[0]{\contentsname}{toc}
\tableofcontents*
\cleardoublepage

% ------------------------------------------------------------------------------
% ELEMENTOS TEXTUAIS
% ------------------------------------------------------------------------------

% Introduz a parte textual
\textual

\chapter{Introdução}

	Parte inicial do texto que visa inserir o leitor no assunto trabalhado na sequência. Deve apresentar de forma geral o trabalho e, em especial, deve abordar o problema, os objetivos e a justificativa do trabalho.

	\section{Problema de pesquisa}
	
		Apresentação do problema que norteará a pesquisa para o TCC. A enunciação do problema deve ser,	preferencialmente, em forma de uma pergunta.
	
	\section{Justificativa}
	
		Texto construído com a intenção de mostrar	a relevância, a importância, a pertinência e a viabilidade do trabalho.
	
	\section{Objetivos}
	
		\subsection{Objetivo geral}
		
			O Objetivo Geral deve ter relação íntima com o problema de	pesquisa e deve apontar o rumo a ser percorrido para encontrar a resposta. Já os	Objetivos Específicos desdobram o Objetivo Geral nos passos necessários para executar o	Objetivo Geral. Os objetivos devem indicar	exatamente a ação a ser tomada.
		
		\subsection{Objetivos específicos}
		
			\begin{enumerate}
				\item Objetivo específico 1;
				\item Objetivo específico 2;
				\item Objetivo específico n.
			\end{enumerate}

\chapter{Revisão da Literatura}
	
	Trata-se de um texto que apresenta de forma	geral os fundamentos teóricos (ou bases	teóricas) e conceituais do trabalho. Deve destacar as principais obras e teorias da área em estudo. Deve estar bem referenciado.
	
	
	\begin{figure*}[h]{
		\caption{\label{Debian}\textit{QRcode Github}}
		\centering
		\includegraphics[width=9cm]{Imagens/Karan_GitHub.png}
		\legend{Fonte: MODELO PARA ADICAO DE IMAGENS}}
	\end{figure*}
	
		\begin{citacao}
		Exemplo de citação direta. \cite{nomenacitacao}
	\end{citacao}
	
	Citação indireta \cite{nomenacitacao}.
	
\chapter{Metodologia}
	
	Deve dizer como o trabalho será realizado. Aborda quatro componentes: descrição do foco do estudo; caracterização da pesquisa (tipo de pesquisa); plano de coleta de dados (técnicas e instrumentos de coleta de dados e informações afins); plano de análise dos dados (técnicas de sistematização e análise dos dados e as formas de apresentação dos resultados).
	
\chapter{Desenvolvimento...}
	
	Demais capítulos do desenvolvimento do trabalho...

\chapter{Considerações finais}

	Parte final do texto na qual são apresentadas as considerações finais acerca dos objetivos, das variáveis e do problema da pesquisa. Não usar esta seção para sumarizar os resultados (o que já foi feito no Resumo), mas destacar o progresso e as aplicações que o trabalho propicia. Enfatizar as limitações que persistem, apresentando, sempre que apropriado, sugestões para trabalhos futuros.
	
\chapter{ORÇAMENTO}
\begin{table}[!h]
	\caption[Orcamento]{Orçamento.}
	\label{tab-nivinv}	
	\begin{tabular}{|c|c|c|}
		\hline
		\textbf{Descrição das Despesas} & \textbf{Quantidade} & \textbf{Valor Estimado (em reais)} \\ \hline
		Fotocópia & 200 & R\$ 100,00 \\ \hline
		Encadernação & 7 & R\$ 50,00 \\ \hline
		Aquisição de bibliografias & 2 & R\$ 60,00 \\ \hline
		Aquisição de cursos & 2 & R\$ 100,00 \\ \hline
		\textbf{Total} &  & \textbf{R\$ 310,00} \\ \hline
	\end{tabular}
\end{table}


% Finaliza o bookmarking do PDF
\phantompart

% ------------------------------------------------------------------------------
% ELEMENTOS PÓS-TEXTUAIS
% ------------------------------------------------------------------------------

% Introduz a parte pós-textual
\postextual

% Insere as referências bibliográficas
\bibliography{bibliografia}

% ----------------- APÊNDICES -----------------
\begin{apendicesenv}

\partapendices

\chapter{Título do Apêndice}

	Documento ou texto “elaborado” pelo autor, inserido ou referido no trabalho a fim de complementar a sua argumentação. Para	não haver prejuízo à unidade nuclear do	trabalho, cópia integral do documento é	acrescida ao final.
	
\chapter{Outro Apêndice}

	Outro documento ou texto “elaborado” pelo autor, inserido ou referido no trabalho a fim de complementar a sua argumentação. Para	não haver prejuízo à unidade nuclear do	trabalho, cópia integral do documento é	acrescida ao final.

% Inserir outros apêndices caso seja necessário ...

\end{apendicesenv}


% ----------------- ANEXOS -----------------
\begin{anexosenv}

\partanexos

\chapter{Título do Anexo}

	Documento ou texto “não elaborado”	pelo autor, inserido ou referido no trabalho, a fim de fundamentar, comprovar ou ilustrar a sua argumentação. Igualmente, para não haver prejuízo à unidade nuclear do trabalho, cópia integral do documento é acrescida ao final do trabalho.
	
\chapter{Outro Anexo}

	Outro documento ou texto “não elaborado”	pelo autor, inserido ou referido no trabalho, a fim de fundamentar, comprovar ou ilustrar a sua argumentação. Igualmente, para não haver prejuízo à unidade nuclear do trabalho, cópia integral do documento é acrescida ao final do trabalho.
	
\end{anexosenv}

\end{document}
